This dissertation presented the design, implementation, testing, deployment, and evaluation of a hybrid medical haematology chatbot built around an explicit MLOps lifecycle. The project addressed a practical gap between general-purpose conversational systems and the narrowly defined needs of haematology laboratory support. General chatbots may provide broad language coverage, but they do not naturally satisfy laboratory requirements for controlled scope, reproducible behaviour, auditability, safe handling of out-of-scope prompts, or continuous operational monitoring. In response to that problem, a domain-specific assistant was developed for English-language haematology workflows, coagulation specimen handling, complete blood count terminology, blood smear support, quality control topics, and report-reading assistance.

The implemented system combined two BiLSTM sequential intent classifiers with retrieval-based response generation, entity-assisted routing, safety rules, PostgreSQL logging, an administrative monitoring panel, and a retraining workflow. A dual-model design was adopted. The first model supported general haematology assistant functions, including workflow and communication intents. The second model targeted blood-report questions, such as CBC parameters, abnormal flags, and report-structure explanations. This design allowed narrower label spaces, targeted monitoring, and model switching when a user query matched the alternative scope more closely.

A fixed 70/15/15 train-validation-test split policy was introduced to strengthen evaluation discipline and reduce leakage between model selection and final reporting. At the time of the final documented evaluation, the master labelled dataset contained 5,391 utterances across 29 intents. The derived general dataset contained 2,670 utterances across 21 intents, and the report dataset contained 4,777 utterances across 26 intents. The general model achieved 0.8928 test accuracy and 0.8649 macro F1, while the report model achieved 0.9358 test accuracy and 0.8965 macro F1 after the report-analysis expansions. In addition, 52 automated tests passed across API, preprocessing, routing, model advisory, split generation, report analysis, and deployment utilities.

It was concluded that the project met its aim of producing a substantial and operational artefact that connected sequential NLP modelling with retrieval-based response generation, report-value information extraction, and practical MLOps controls. The main strengths of the work were controlled deployment scope, reproducible data handling, traceable inference, bounded report analysis, and maintainable operational monitoring. The main limitations were the use of a moderate-sized labelled dataset, English-only support, the absence of patient-specific clinical interpretation, and continued reliance on rule-based entity refinement and reference logic. Nevertheless, the project demonstrated that a hybrid sequential architecture remained appropriate for a bounded medical assistant when supported by disciplined data management, monitoring, and retraining processes.
